\documentclass[10pt]{article}
\usepackage[a4paper,landscape,margin=1.4cm]{geometry}
\usepackage{graphicx}
\usepackage{booktabs}
\usepackage{enumitem}
\usepackage{siunitx}
\usepackage[hidelinks]{hyperref}
\usepackage{parskip}
\setlist{itemsep=2pt,topsep=3pt}

\newcommand{\atlas}[1]{\begin{center}\includegraphics[width=\linewidth,height=0.74\textheight,keepaspectratio]{figures/#1}\end{center}}

\title{Passive cardiac shear-wave elastography:\\ processing methods explored, as space-time plots}
\author{shearWaveProcessing -- \texttt{study/analysis/passive\_methods\_atlas.py}}
\date{24 September 2026}

\begin{document}
\maketitle

\section*{What this document is}

The passive (valve-closure) shear-wave work compared many processing choices: which motion
quantity to show, the temporal band-pass, spatial and temporal smoothing, M-line averaging,
directional filtering, clutter and motion handling, and recipes from the literature. The
decisions were recorded in the repo's docs (\texttt{docs/passive\_search.md}, \texttt{HANDOFF.md}), but mostly
on one acquisition and often through automatic scores. This document shows the choices
\emph{side by side on the same in-vivo windows}, as space-time (M-mode) plots, so they can be
judged by eye. \textbf{No automatically fitted speed is shown anywhere}; it is not what we trust
(\texttt{docs/passive\_speed\_estimation.md}).

All data are in vivo (buffer 4, the $\sim$\SI{926}{\hertz} diverging-wave stream of the S5-1
sequence). The phantom has no passive shear waves.

\subsection*{How to read the panels}
\begin{itemize}
  \item \textbf{Axes.} Horizontal: time relative to the detected event peak, in ms (dotted line
    at 0). Vertical: distance along the analysed part of the M-line, in mm. So a wave travelling
    along the septum is a \emph{tilted} band, and motion of the whole wall at once is a
    \emph{vertical} band.
  \item \textbf{How much tilt to expect.} A wave at \SI{3}{\metre\per\second} crosses a \SI{20}{\milli\metre}
    line in about \SI{7}{\milli\second}. That is small on a \SI{130}{\milli\second} axis, so
    Figure~\ref{fig:zoom} repeats the key comparison zoomed to $\pm$\SI{25}{\milli\second}.
  \item \textbf{Colour.} Each panel has its own symmetric colour scale (99th percentile of
    $|\cdot|$), so the plots show pattern, not amplitude.
  \item \textbf{Columns.} Seven windows, always in the same order. Six were scored \emph{clear} in
    the hand-labelled set: three aortic valve closures (AVC) and three mitral valve closures (MVC).
    The last column was scored \emph{none}: it shows what processing makes of a window with no
    visible wave. \textbf{A method that makes the last column look like a wave is manufacturing structure.}
  \item \textbf{Timing.} The detected event peak and the visible wave do not always coincide:
    in C000000003 the wave is $\sim$\SI{40}{\milli\second} before the peak.
\end{itemize}

\subsection*{Baseline recipe}
Every family starts from view A of \texttt{configs/passive.yaml} and changes one thing. View A:
\begin{itemize}
  \item Loupas lag-1 estimator, frame-to-frame, cumulative displacement;
  \item 10--150\,Hz zero-phase Butterworth band-pass;
  \item Gaussian spatial smoothing, $\sigma$ = 0.6\,mm axial $\times$ 1.2\,mm lateral;
  \item 3-frame moving mean;
  \item 5 parallel M-lines 0.5\,mm apart, averaged;
  \item no directional filter.
\end{itemize}
The exceptions are the literature recipes and the three production views, which are complete
recipes. The IQ is cropped to the M-line box plus 8\,mm before processing.

\begin{figure}[h]
\atlas{windows.png}
\caption{The seven windows: the buffer-1 B-mode frame at the event's cardiac phase, on which the
per-event M-line was drawn. Yellow: the part of the line analysed (the part hand-labelled);
grey: the rest of the line. M-lines are 19--22\,mm.}
\end{figure}

\clearpage
\section{Motion quantity: displacement, velocity, acceleration}

\textbf{What varies.} The same frame-to-frame estimate shown as cumulative displacement, as
particle velocity, or as acceleration (the temporal derivatives).

\textbf{Literature.} Velocity is used by Keijzer et al.\ (2019, 2020). Acceleration is used by
Petrescu et al., Santos et al.\ (2019), Strachinaru et al.\ and Espeland et al.\ (2024).
Nobody reports cumulative displacement. The wave is dispersive, so the three weight different
frequencies and give different apparent speeds.

\textbf{Earlier conclusion.} Displacement was chosen for visualisation (cleanest), velocity
was noisier, and acceleration had the sharpest front but was noisiest.

\textbf{What the atlas shows.}
\begin{itemize}
  \item \textbf{Displacement}: smooth but broad bands, roughly \SIrange{15}{25}{\milli\second}
    wide. At this width a tilt of a few ms is hard to read.
  \item \textbf{Velocity}: a bipolar front about half as wide. The tilt becomes visible in
    C000000019, 020 and 002 (Figure~\ref{fig:zoom}).
  \item \textbf{Acceleration}: the narrowest front exactly at the event (019, 020, 027 -- in 027
    the clearest tilted line of all). But outside the event it is noise-dominated, and in 003 and
    039 the front is lost.
  \item \textbf{The no-wave window} is incoherent in velocity and acceleration. Displacement turns it
    into smooth broad blobs that are easier to misread.
\end{itemize}

\atlas{quantity.png}

\begin{figure}[p]
\atlas{quantity_zoom.png}
\caption{The quantity comparison of the previous page zoomed to $\pm$\SI{25}{\milli\second} around the midpoint of the hand-drawn
wavefront (event peak for the no-wave window). Tilted fronts are readable in velocity (019, 020,
002, 003) and in acceleration (027). Displacement bands are too wide to show the slope well.}
\label{fig:zoom}
\end{figure}

\clearpage
\section{Temporal band-pass}

\textbf{What varies.} The zero-phase Butterworth band-pass that separates the wave from bulk wall
motion (low corner) and from noise (high corner).
\begin{itemize}
  \item 15--100\,Hz is Keijzer's band and the Caenen review's band for MVC waves.
  \item 15--90\,Hz is view C.
  \item 8--45\,Hz is the narrowing tried on 2026-09-22 and reverted.
  \item 30--150\,Hz tests a higher low corner.
\end{itemize}

\textbf{Earlier conclusion.} The low corner must be at least 10--20\,Hz, because 5\,Hz lets the
flat bulk band through. The band content ends near 90\,Hz (f99), and 8--45\,Hz cut real signal.

\textbf{What the atlas shows (displacement first, then velocity).}
\begin{itemize}
  \item \textbf{5--150\,Hz} leaves large slow blobs (e.g.\ 019 at $-20$\,ms, 020 and 027 broad):
    bulk wall motion.
  \item \textbf{10--150\,Hz} removes most of it.
  \item \textbf{15--100 and 15--90\,Hz} give the narrowest clean fronts in both quantities. In
    velocity they are the cleanest panels of the whole atlas, and the no-wave window stays
    incoherent.
  \item \textbf{8--45\,Hz} smears every front into a broad band.
  \item \textbf{30--150\,Hz} turns each event into several narrow parallel bands (ringing): the low
    corner now sits inside the wave's own spectrum, and the no-wave window becomes fine texture.
\end{itemize}

\atlas{band_displacement.png}
\atlas{band_velocity.png}

\clearpage
\section{Directional (k--$\omega$) filter}

\textbf{What varies.} The filter keeps only waves travelling one way along the line (toward low
or toward high $r$). It is standard for ARF shear waves (Song et al.\ 2016). For natural waves
none of the papers use one.

\textbf{Earlier conclusion.} Keep it off for passive: it added banding and biased the slope. The
synthetic tests of 2026-09-24 also show it biases the active speed high.

\textbf{What the atlas shows.} \textbf{It imposes a tilt on everything.}
\begin{itemize}
  \item In the no-wave window it produces clean, regular tilted stripes -- a convincing-looking
    wave that is not there.
  \item In the real windows the apparent tilt follows the direction chosen: the same event tilts
    one way with ``toward low $r$'' and the other way with ``toward high $r$'' (019, 020, 039).
  \item The filter therefore decides the direction instead of measuring it. It must stay off
    whenever the direction is not known in advance.
\end{itemize}

\atlas{directional.png}

\clearpage
\section{Spatial smoothing}

\textbf{What varies.} Per-frame smoothing of the displacement field before sampling the M-line.

\textbf{Earlier conclusion.} Light Gaussian (0.6--1.0\,mm) is best; median is similar or
slightly worse.

\textbf{What the atlas shows.}
\begin{itemize}
  \item \textbf{None}: fine horizontal stripes along $r$ (speckle decorrelation between depths).
  \item \textbf{Gaussian 0.6\,$\times$\,1.2\,mm} removes them.
  \item \textbf{Gaussian 1.0\,$\times$\,2.0\,mm} changes little beyond that.
  \item \textbf{Gaussian 2\,$\times$\,4\,mm} makes the fronts look clean. But it turns the no-wave
    window into smooth coherent bands as well: too much.
  \item \textbf{Median 1\,$\times$\,2\,mm} sits between none and Gaussian and keeps some stripes.
\end{itemize}

\atlas{spatial.png}

\clearpage
\section{Temporal smoothing}

\textbf{What varies.} A moving mean or median over 3--5 frames (3--5\,ms).

\textbf{Earlier conclusion.} Second-order.

\textbf{What the atlas shows.} Almost no visible difference between none, mean 3, mean 5 and
median 5. The 150\,Hz low-pass corner already limits the bandwidth more than these windows do.
Temporal smoothing matters for velocity and acceleration (derivatives amplify noise), not for
band-passed displacement.

\atlas{temporal.png}

\clearpage
\section{M-line averaging}

\textbf{What varies.} How many parallel copies of the M-line, offset along its normal, are
averaged.

\textbf{Earlier conclusion.} Few lines. One line is best for the sharp acceleration front; more
lines help smoother quantities.

\textbf{What the atlas shows.}
\begin{itemize}
  \item Little change from 1 to 9 lines in displacement.
  \item At 15 lines (a $\pm$5.6\,mm band -- more than half the wall thickness) the no-wave window
    becomes smoother and more band-like, and near the wall edges (003, 002) the pattern changes.
\end{itemize}

\atlas{mline.png}

\clearpage
\section{Motion and clutter handling}

\textbf{What varies.} What removes the bulk wall motion:
\begin{itemize}
  \item the band-pass alone (view A);
  \item plus an IQ slow-time low-pass at 250\,Hz (Keijzer's blood/noise filter);
  \item plus an SVD clutter filter on the IQ (one component, Demen\'e-style);
  \item a polynomial detrend instead of the band-pass high side;
  \item or only a 5\,Hz high-pass, which leaves the bulk motion in as a reference.
\end{itemize}

\textbf{Earlier conclusion.} The low corner of the band-pass is the decisive motion remover. SVD
and polynomial detrending do not substitute for it, and the 250\,Hz IQ low-pass washed out the
wave in the search.

\textbf{What the atlas shows.}
\begin{itemize}
  \item With the 10\,Hz high-pass in place, the IQ low-pass and the SVD filter leave the panels
    essentially unchanged.
  \item The polynomial detrend adds bands at both window ends (the polynomial does not fit the
    edges) and keeps more low-frequency content.
  \item The 5\,Hz-only row is dominated by broad bulk motion.
\end{itemize}

\atlas{motion.png}

\clearpage
\section{Recipes from the literature}

\begin{itemize}
  \item \textbf{Keijzer et al.} IQ slow-time low-pass at 250\,Hz; lag-1 velocity with Gaussian
    smoothing of the autocorrelation before the angle (4\,mm axial; their lateral 6.7 is
    approximated as 3\,mm); 15--100\,Hz band-pass (6th order).
  \item \textbf{Petrescu / Santos.} Acceleration, with a 3-frame moving mean on the velocity
    before differentiation.
  \item \textbf{Espeland et al.\ 2024 (clutter filter wave imaging, CFWI, PLAX view).} IQ
    high-pass at a 2\,cm/s cutoff, envelope, temporal derivative. The first and last 10 frames
    are dropped (filter end transient).
\end{itemize}

\textbf{What the atlas shows.}
\begin{itemize}
  \item \textbf{Keijzer's velocity recipe} looks like the 15--100\,Hz velocity rows: clean bipolar
    fronts, and an incoherent no-wave window.
  \item \textbf{The acceleration recipe} has sharp fronts at 019, 020 and 027 and is
    noise-dominated elsewhere.
  \item \textbf{CFWI} shows the event as a front in 002, 027 and 039, faintly in 019 and 020, with
    lower contrast and vertical streaks. Its automatic direction agreed with the hand picks only at
    chance level (\texttt{study/analysis/cfwi\_benchmark.py}).
\end{itemize}

\atlas{literature.png}

\clearpage
\section{The three production views}

\texttt{configs/passive.yaml} renders three deliberately different recipes per window:
\begin{itemize}
  \item \textbf{A}: displacement, 10--150\,Hz;
  \item \textbf{B}: displacement, 5--150\,Hz, median 1\,$\times$\,2\,mm, no temporal smoothing,
    9 lines;
  \item \textbf{C}: velocity, 15--90\,Hz, Gaussian 1.0\,$\times$\,2.0\,mm, mean 5, 9 lines.
\end{itemize}
In these windows:
\begin{itemize}
  \item \textbf{C} gives the sharpest, most readable fronts and the most clearly incoherent
    no-wave window.
  \item \textbf{B} carries bulk-motion blobs from its 5\,Hz corner (019, 027) and median streaks.
  \item \textbf{A} is in between: smooth but broad.
\end{itemize}

\atlas{views.png}

\clearpage
\section{Summary}

\begin{center}
\small
\begin{tabular}{p{4.2cm}p{11.2cm}p{7.2cm}}
\toprule
choice & what the space-time plots show (7 windows) & suggestion for visual reading \\
\midrule
quantity & displacement smooth but broad (slope hard to read); velocity half as wide, tilt visible;
  acceleration sharpest at the event, noisy elsewhere & read slopes on \textbf{velocity}; use
  acceleration as a timing cross-check; displacement for robustness only \\
band-pass & 5\,Hz keeps bulk motion; 15--90/100\,Hz narrowest clean fronts; 8--45 smears;
  30--150 rings & \textbf{15--100\,Hz} (literature band) \\
directional filter & imposes a tilt, including on the no-wave window; direction follows the
  setting & \textbf{off} \\
spatial smoothing & none striped; 0.6--1.0\,mm Gaussian clean; 2\,$\times$\,4\,mm fabricates bands
  in noise & Gaussian \textbf{0.6--1.0\,mm} \\
temporal smoothing & no visible effect on band-passed displacement & mean 3--5 on velocity /
  acceleration \\
M-line averaging & little effect up to 9 lines; 15 lines smooths noise into bands & \textbf{5--9 lines}
  within the wall \\
IQ low-pass / SVD & no visible change once the high-pass is in place & not needed \\
polynomial detrend & edge bands, more low-frequency content & not a substitute for the band-pass \\
CFWI & fronts in 3--5 of 6 windows at lower contrast; direction at chance & not better here \\
\bottomrule
\end{tabular}
\end{center}

\textbf{Caveats.}
\begin{itemize}
  \item Seven windows, labelled by one observer; the confidence scores were not blind.
  \item The lines are short (19--22\,mm) compared with the shear wavelength (150\,mm or more at
    20\,Hz). Much of the ``front'' is therefore the whole segment moving nearly together, and a
    method can look good by making bands vertical.
  \item Per-panel colour scaling hides amplitude differences.
  \item These observations favour a velocity view in the $\sim$15--100\,Hz band (close to view C
    and to Keijzer's recipe) over displacement for reading slopes. That differs from the earlier
    ``displacement is cleanest'' conclusion, which was reached on one acquisition. A second
    observer's reading (\texttt{docs/TODO.md}, item 4) should settle it.
\end{itemize}

\textbf{Reproduce.}
\begin{itemize}
  \item \texttt{python study/analysis/passive\_methods\_atlas.py} reads the seven windows from
    \texttt{Z:/raw\_data} ($\sim$20\,min) and writes \texttt{report/passive\_methods/figures/}.
  \item \texttt{--figures} re-renders from the cache.
  \item Compile with \texttt{tectonic passive\_methods.tex}.
\end{itemize}

\end{document}
